\documentclass[11pt,a4paper]{altacv}

\geometry{left=1.5cm,right=1.5cm,top=1.cm,bottom=1.2cm,columnsep=1.5cm}

\usepackage[utf8]{inputenc}
\usepackage[T1]{fontenc}
\usepackage[french]{babel}
\usepackage{paracol}
\usepackage{xcolor}
\usepackage{hyperref}
\usepackage{fontawesome5}
\usepackage{setspace}

\definecolor{SlateGrey}{HTML}{2E2E2E}
\definecolor{LightGrey}{HTML}{666666}
\definecolor{DarkPastelRed}{HTML}{8AC0D9}
\definecolor{PastelRed}{HTML}{00B0DA}
\definecolor{GoldenEarth}{HTML}{B4AD0C}
\colorlet{name}{black}
\colorlet{tagline}{PastelRed}
\colorlet{heading}{DarkPastelRed}
\colorlet{headingrule}{GoldenEarth}
\colorlet{subheading}{PastelRed}
\colorlet{accent}{PastelRed}
\colorlet{emphasis}{SlateGrey}
\colorlet{body}{LightGrey}

\renewcommand{\familydefault}{\sfdefault}

\name{\Large Guillaume BOILEAU}
\tagline{Consultant Analytique Energie / Data Scientist Industriel | ML, IoT, performance, optimisation \& reporting}
\personalinfo{
  \email{guillaume.boileaupro@gmail.com}
  \phone{+33 6 59 94 85 84}
  \location{Alpes-Maritimes, France}
  \linkedin{guillaume-boileau-phd-891b81265}
}

\setlength{\parskip}{0.72\baselineskip}

\begin{document}
\makecvheader
\vspace{-0.65cm}

\cvsection{Lettre de motivation}
\vspace{-0.45cm}

{\color{accent}\textbf{Objet :}} \textit{Candidature -- Consultant Analytique Energie -- Carros}

{\color{body}

Madame, Monsieur,

Je vous adresse ma candidature pour le poste de \textbf{Consultant Analytique Energie}. Votre offre m'intéresse particulièrement car elle relie plusieurs dimensions qui correspondent fortement à mon parcours : analyse de données opérationnelles, machine learning, données IoT, optimisation de performance, validation de modèles, reporting et recommandations concrètes pour des équipes métiers et clients.

Docteur en physique et ingénieur R\&D / data, j'ai construit mon parcours autour de l'analyse de systèmes complexes, du développement de pipelines Python, de la modélisation statistique, de la validation de résultats et de la transformation de données techniques en éléments exploitables pour la décision. Ce qui m'attire dans ce poste est précisément cette logique : partir de données terrain, évaluer leur qualité, identifier des opportunités d'optimisation, tester des modèles, suivre leur comportement dans le temps et restituer des recommandations utiles.

Mon expérience récente chez \textbf{VEV Platform Services France} constitue un lien concret avec les sujets énergie et infrastructures connectées. J'y ai travaillé sur une plateforme logicielle CPMS liée à des infrastructures de recharge de véhicules électriques. Cette expérience m'a confronté à des problématiques très opérationnelles : flux de données, incidents, incohérences, contraintes terrain, fiabilité des échanges et analyse des comportements applicatifs. Elle m'a appris à relier données, infrastructure, exploitation et qualité de service dans un environnement où les résultats doivent être compréhensibles et actionnables.

Au \textbf{CNES}, j'ai travaillé sur le segment sol de la mission spatiale LISA, dans un environnement CNES/ESA associant données complexes, simulation, intégration de modèles et validation technique. J'y ai développé \textbf{CATpyLISA}, une plateforme Python destinée à structurer des workflows, comparer des modèles, contrôler la cohérence des résultats et documenter les écarts. Cette expérience m'a donné une pratique solide de la construction de pipelines reproductibles, de la validation de modèles, de la qualité de données et du reporting technique.

À l'Universiteit Antwerpen, j'ai également travaillé sur des données capteurs, des séries temporelles, des perturbations environnementales et des indicateurs de performance pour instruments de précision. J'y ai développé des workflows de data science et de machine learning pour analyser des mesures, réduire des bruits non stationnaires, comparer des configurations et produire des recommandations techniques. Cette expérience est directement transposable à l'analyse de données IoT et à l'identification d'opportunités d'optimisation à partir de signaux opérationnels.

Pour Schneider Electric, je peux apporter un profil capable de comprendre des données industrielles, d'évaluer leur qualité, de construire ou tester des modèles, de suivre leurs performances et de formuler des recommandations claires. Je suis à l'aise avec Python, pandas, NumPy, SciPy, scikit-learn, SQL, Power BI, Linux, Git, Docker et les pratiques de documentation. Je peux également travailler en interface avec des experts énergie, des équipes projet, des développeurs, des responsables client et des utilisateurs métiers.

Je suis conscient que certains éléments spécifiques du poste demandent une montée en compétence ciblée, notamment les plateformes ML low-code/no-code, l'écosystème AVEVA/Schneider, ainsi que les systèmes industriels de type eau glacée, air comprimé, pompage, ventilation, vapeur ou traitement d'air. En revanche, mon expérience m'a donné une capacité solide à entrer rapidement dans un domaine technique, à comprendre ses données, ses contraintes et ses indicateurs, puis à produire des analyses robustes et utiles.

Ce poste représente pour moi une évolution cohérente vers un rôle à impact, au croisement de la data science, de l'énergie, de l'industrie et de la durabilité. Contribuer à optimiser la consommation énergétique, à améliorer la performance des installations et à transformer des cas d'usage ML en bibliothèques réutilisables correspond très directement à ce que je souhaite faire aujourd'hui : mettre mes compétences quantitatives et logicielles au service de solutions concrètes.

Je serais heureux de pouvoir échanger avec vous afin de vous présenter plus en détail mon parcours et la manière dont mon expérience en data science, machine learning, données capteurs, validation de modèles, reporting et mobilité électrique peut contribuer aux projets d'optimisation énergétique portés par Schneider Electric et AVEVA.

Je vous prie d'agréer, Madame, Monsieur, l'expression de mes salutations distinguées.

\textbf{Guillaume Boileau}

}

\end{document}
